% Source: https://artofproblemsolving.com/community/c3h1753532_counting_triangles
% Stable identity: aops-topic-1753532-q01
\textbf{(a)} No three vertices of a regular hexagon are collinear, and a triangle is determined by its three vertices, so there are $\dbinom63=\boxed{20}$ triangles.

\textbf{(b)} Classify the $20$ triangles by the arcs between their vertices around the hexagon (the three gaps, measured in hexagon steps, add up to $6$): the gap patterns are $\{2,2,2\}$, $\{1,1,4\}$ and $\{1,2,3\}$.

\begin{asy}
usepackage("amsmath");
size(9cm);
defaultpen(fontsize(10pt));
pair[] P; for(int i=0;i<6;++i) P[i]=dir(60*i);
picture hexp(int k){picture pic; pair[] Q;
 for(int i=0;i<6;++i) Q[i]=P[i];
 filldraw(pic,(k==0?Q[0]--Q[2]--Q[4]--cycle:(k==1?Q[0]--Q[1]--Q[2]--cycle:Q[0]--Q[1]--Q[3]--cycle)),paleblue,blue+linewidth(0.8));
 draw(pic,Q[0]--Q[1]--Q[2]--Q[3]--Q[4]--Q[5]--cycle,linewidth(0.8));
 for(int i=0;i<6;++i) dot(pic,Q[i]);
 return pic;}
add(shift(0,0)*hexp(0));add(shift(2.6,0)*hexp(1));add(shift(5.2,0)*hexp(2));
label("equilateral",(0,-1.3));label("isosceles $120^\circ$",(2.6,-1.3));label("$30^\circ$-$60^\circ$-$90^\circ$",(5.2,-1.3));
label("\mbox{}",(min(currentpicture,user=true)+max(currentpicture,user=true))/2,black); // keeps the text after the figure black when LaTeX typesets the labels inline
\end{asy}

\begin{itemize}
\item $\{2,2,2\}$: alternate vertices, an equilateral triangle with side $\sqrt3$ (for unit side): $2$ triangles.
\item $\{1,1,4\}$: three consecutive vertices, an isosceles triangle with sides $1,1,\sqrt3$ and a $120^\circ$ angle: $6$ triangles.
\item $\{1,2,3\}$: two opposite vertices and one more, a $30^\circ$--$60^\circ$--$90^\circ$ triangle with sides $1,\sqrt3,2$: $12$ triangles.
\end{itemize}
Check: $2+6+12=20$. Triangles in the same class are congruent (a rotation or reflection of the hexagon maps one to another), and the three classes have different side lengths. So there are $\boxed{3}$ non-congruent triangles (not $4$).
